%File: formatting-instructions-latex-2024.tex
%release 2024.0
\documentclass[letterpaper]{article} % DO NOT CHANGE THIS
\usepackage{aaai24}  % DO NOT CHANGE THIS
\usepackage{times}  % DO NOT CHANGE THIS
\usepackage{helvet}  % DO NOT CHANGE THIS
\usepackage{courier}  % DO NOT CHANGE THIS
\usepackage[hyphens]{url}  % DO NOT CHANGE THIS
\usepackage{graphicx} % DO NOT CHANGE THIS
\urlstyle{rm} % DO NOT CHANGE THIS
\def\UrlFont{\rm}  % DO NOT CHANGE THIS
\usepackage{natbib}  % DO NOT CHANGE THIS AND DO NOT ADD ANY OPTIONS TO IT
\usepackage{caption} % DO NOT CHANGE THIS AND DO NOT ADD ANY OPTIONS TO IT
\frenchspacing  % DO NOT CHANGE THIS
\setlength{\pdfpagewidth}{8.5in}  % DO NOT CHANGE THIS
\setlength{\pdfpageheight}{11in}  % DO NOT CHANGE THIS
%
% These are recommended to typeset algorithms but not required. See the subsubsection on algorithms. Remove them if you don't have algorithms in your paper.
\usepackage{algorithm}
\usepackage{algorithmic}
\usepackage{booktabs}  % For professional-quality tables
%
% These are are recommended to typeset listings but not required. See the subsubsection on listing. Remove this block if you don't have listings in your paper.
\usepackage{newfloat}
\usepackage{listings}
\DeclareCaptionStyle{ruled}{labelfont=normalfont,labelsep=colon,strut=off} % DO NOT CHANGE THIS
\lstset{%
	basicstyle={\footnotesize\ttfamily},% footnotesize acceptable for monospace
	numbers=left,numberstyle=\footnotesize,xleftmargin=2em,% show line numbers, remove this entire line if you don't want the numbers.
	aboveskip=0pt,belowskip=0pt,%
	showstringspaces=false,tabsize=2,breaklines=true}
\floatstyle{ruled}
\newfloat{listing}{tb}{lst}{}
\floatname{listing}{Listing}
%
% Keep the \pdfinfo as shown here. There's no need
% for you to add the /Title and /Author tags.
\pdfinfo{
/TemplateVersion (2024.1)
}

% DISALLOWED PACKAGES
% \usepackage{authblk} -- This package is specifically forbidden
% \usepackage{balance} -- This package is specifically forbidden
% \usepackage{color (if used in text)
% \usepackage{CJK} -- This package is specifically forbidden
% \usepackage{float} -- This package is specifically forbidden
% \usepackage{flushend} -- This package is specifically forbidden
% \usepackage{fontenc} -- This package is specifically forbidden
% \usepackage{fullpage} -- This package is specifically forbidden
% \usepackage{geometry} -- This package is specifically forbidden
% \usepackage{grffile} -- This package is specifically forbidden
% \usepackage{hyperref} -- This package is specifically forbidden
% \usepackage{navigator} -- This package is specifically forbidden
% (or any other package that embeds links such as navigator or hyperref)
% \indentfirst} -- This package is specifically forbidden
% \layout} -- This package is specifically forbidden
% \multicol} -- This package is specifically forbidden
% \nameref} -- This package is specifically forbidden
% \usepackage{savetrees} -- This package is specifically forbidden
% \usepackage{setspace} -- This package is specifically forbidden
% \usepackage{stfloats} -- This package is specifically forbidden
% \usepackage{tabu} -- This package is specifically forbidden
% \usepackage{titlesec} -- This package is specifically forbidden
% \usepackage{tocbibind} -- This package is specifically forbidden
% \usepackage{ulem} -- This package is specifically forbidden
% \usepackage{wrapfig} -- This package is specifically forbidden
% DISALLOWED COMMANDS
% \nocopyright -- Your paper will not be published if you use this command
% \addtolength -- This command may not be used
% \balance -- This command may not be used
% \baselinestretch -- Your paper will not be published if you use this command
% \clearpage -- No page breaks of any kind may be used for the final version of your paper
% \columnsep -- This command may not be used
% \newpage -- No page breaks of any kind may be used for the final version of your paper
% \pagebreak -- No page breaks of any kind may be used for the final version of your paperr
% \pagestyle -- This command may not be used
% \tiny -- This is not an acceptable font size.
% \vspace{- -- No negative value may be used in proximity of a caption, figure, table, section, subsection, subsubsection, or reference
% \vskip{- -- No negative value may be used to alter spacing above or below a caption, figure, table, section, subsection, subsubsection, or reference

\setcounter{secnumdepth}{0} %May be changed to 1 or 2 if section numbers are desired.

% The file aaai24.sty is the style file for AAAI Press
% proceedings, working notes, and technical reports.
%

\title{Machine Learning for Exoplanet Detection: Identifying Exoplanets in Light Curves}
\author {
    % Authors
    Tristan Moffett\textsuperscript{\rm },
    Josh Manchester\textsuperscript{\rm },
    Brianne Leatherman\textsuperscript{\rm }
}
\affiliations {
    % Affiliations
    \textsuperscript{\rm }University of Colorado at Colorado Springs\\
    %\textsuperscript{\rm 2}Affiliation 2\\
    %firstAuthor@affiliation1.com, secondAuthor@affilation2.com, thirdAuthor@affiliation1.com
}

\begin{document}
\maketitle

\section*{Abstract}
The discovery of exoplanets has become more dependent on machine learning as missions like Kepler and TESS generate thousands of transit signals that are too costly and time consuming to vet by hand. In this research, we address this problem by comparing three different deep learning models for exoplanet classification: Convolutional Neural Networks (CNNs), Recurrent Neural Networks (RNNs), and Transformer architecture. Each team member designed, trained, and evaluated one of these models using observational data from Kepler and the Transiting Exoplanet Survey Satellite (TESS) along with simulated datasets. Our focus is on performance metrics such as accuracy, precision, recall, F1 score, and Area Under Curve (AUC), as well as understanding the strengths and weaknesses of each model in classifying planets versus false positives.

This paper presents final results from the BiLSTM + clustering approach, which achieved an AUC of 0.9261 on held-out test data (6,579 windows from TESS Sector 1 ground truth) with 100\% recall---detecting every planet in the test set with zero false negatives. The RNN implementation combines K-means clustering on statistical features with a 4-layer bidirectional LSTM architecture, enabling the network to learn cluster-specific patterns and distinguish different types of stellar variability from genuine planetary transits. The model uses class weighting (pos\_weight = 7.41) to handle the 11.9\% positive class imbalance and achieves strong predictive power on the 33,051-window dataset. From analyzing results across the three architectures, we aim to identify which methods are most effective for generalizability across multiple datasets and advance automated exoplanet vetting for future missions.

\section*{Introduction \& Background}
In the last twenty years, there have been significant increases in the number of space missions whose goal is to observe distant stars and the changes that may occur over variable lengths of time. NASA has been one of the primary drivers behind these advances, which are called transit survey missions. Of most importance to our research are two missions in particular: Kepler/K2 and TESS. Kepler was launched in 2009 with the original mission of monitoring a 116 square degree segment of sky for three and a half years, and when the satellite had a mechanical issue, the mission was changed to K2, which monitored stars and systems along the elliptic plane until it was retired in 2018 \cite{salinas2023distinguishingtransitfalsepositives}. The Transiting Exoplanet Survey Satellite (TESS) was launched in 2018, nearly ten years after Kepler, with the goal of monitoring a sector of approximately 2300 square degrees for about 27 days before moving to the next section \cite{salinas2023distinguishingtransitfalsepositives}. At the time of this writing, these two transit surveys have accounted for more than 3000 of the 6000 total confirmed exoplanets \cite{salinas2025tessfullframe}.

The most prevalent method of determining the presence of an extra-orbital planet has been through observing changes in the intensity of the brightness of a star when an object transits between the star and the observer. These transits can be found by representing the observed flux from the star in a time-series graph, also known as light curves, and then analyzing the data \cite{morvan2022dontpayattentionnoise}. This process often includes phase-folding the data by overlapping periods to determine patterns \cite{salinas2023distinguishingtransitfalsepositives}. A dip in the light curve indicates a threshold crossing event (TCE) but that alone does not mean a planet caused the obstruction of the light. The TCE can also be a false positive that was triggered by an eclipsing binary or instrumental artifact \cite{salinas2025tessfullframe}.

The similar nature of the light curve of a transiting exoplanet and that of an eclipsing binary can be difficult to differentiate. Additionally, light curves can be challenging to work with due to the continuous nature of background noise from instrumental, photon, and other interference in the transit of light, and the curves can require significant work to pre-process and then analyze the data \cite{morvan2022dontpayattentionnoise}. As deep learning continues to develop, the use of neural networks has become a common tool in this process, significantly reducing the number of light curves that need to be analyzed by hand.

The Convolutional Neural Network (CNN) was introduced as a model designed to detect repeating patterns in data. For exoplanet research, this means recognizing the U-shaped dips of planetary transits and separating them from the V-shaped dips of eclipsing binaries within time-series light curve data. A challenge for CNNs is that light curves often contain a large amount of noise, which can make training more difficult. However, when handled properly, CNNs can still learn to use this noise to improve classification. \cite{Dattilo2019identifyexoplanetsII} demonstrated the value of CNNs on Kepler data by combining global and local views of a light curve to capture both the orbital view and the details of a transit. \cite{Osborn2019rapidclassificationtesscandidates} later applied a similar approach to TESS data and introduced additional features such as centroid and stellar parameters to handle noisier signals. These results show why CNNs are a strong fit for exoplanet classification, though other deep learning approaches like Recurrent Neural Networks (RNNs) and Transformers bring their own strengths for working with time-series data such as light curves.

Recurrent Neural Networks (RNNs) are a natural fit for time-series data because they learn patterns across many time steps and keep context memory (history) from earlier points in the sequence. In light-curve analysis, this lets the model use both the local transit shape of the light curve and the surrounding temporal pattern (periodic repeats) at the same time. Gated variants like LSTM and GRU help with long-range dependencies and noisy sequences of data. Stacked LSTMs have shown to be a strong resource for space-mission photometry: for example, an LSTM detector trained on Kepler flares generalized well to TESS with high precision/recall when preprocessing and class balancing were handled carefully \cite{Vida2021findingflaresrecurrent}. Beyond direct classification, sequence-aware representation learning with an Echo State Network (ESN) + autoencoder reconstructs light curves in sequence space (not just a readout vector), producing more faithful time-series embeddings for future tasks \cite{kugler2016explorativeapproachkepler}. Finally, RMTPP shows how an RNN can encode event history to predict both the timing and type of the next event; adding a lightweight timing cue (intervals/slight inaccurate period) alongside the classifier can reduce aperiodic false positives in transit searches \cite{Du2016markedtemporalpoint}. In this project, we will benchmark LSTM/GRU on our curated subset, report precision, recall, F1, and ROC-AUC, and optionally include a small timing helper to capture periodicity.

The Transformer architecture was introduced as a novel, simple network relying on encoder/decoder layers of self-attention mechanisms \cite{vaswani2017attentionneed}. Because attention can be calculated in parallel, this network saw an improvement in quality of results with a reduced amount of training time compared to prior RNN and CNN approaches \cite{vaswani2017attentionneed}. Although Transformers were originally designed for Natural Language Processing (NLP), their use in other fields has rapidly spread, due to both the efficiency and the ability to correlate sequential data like time series such as a light curve, where the transit of a planet can be considered a function of time \cite{salinas2023distinguishingtransitfalsepositives}. Given the rapid advancements in deep learning methods over the last decade, a direct performance comparison across RNNs, CNNs, and Transformers on light curve analysis should prove interesting. This paper will detail the methods for creating such a comparison, as well as results, by modeling multiple research processes from other papers.

\section*{Related Work}
\subsection{Convolutional Neural Networks (CNNs)}
Shallue and Vanderburg (\citeyear{Shallue2018identifyexoplanetsI}) addressed the problem of classifying Threshold Crossing Events (TCEs) from the Kepler mission, where thousands of signals had to be examined and most turned out not to be planets. To resolve this bottleneck, they designed one of the first convolutional neural networks (CNNs) for exoplanet detection, comparing a linear model, a fully connected network, and a 1-D CNN, with the CNN proving most successful. Their final architecture called AstroNet processed two versions of each light curve: a global view covering the full orbital phase and a local view zoomed in around the transit. Each view passed through layers of one-dimensional convolutions and max pooling before being merged into fully connected layers of 512 nodes, with dropout and data augmentation used to reduce overfitting. On Kepler's real observational test set the CNN achieved results with 96\% accuracy, 98.8\% AUC, 96\% precision, and 95\% recall. However, when evaluated on simulated datasets the model struggled, recovering 81\% of true injected planets but incorrectly labeling around 90\% of background planets and eclipsing binaries as planets, misclassifying about 65\% of eclipsing binaries, and letting through 3.18\% of false positives from inverted light curves. These outcomes demonstrated that while the model was powerful on Kepler data, it lacked generalizability, leading the authors to suggest training with simulated data and including centroid information to better handle false positives. This highlights the need for models that can work across both real and simulated datasets, a direction we plan to take in our own work.

Dattilo et al. (\citeyear{Dattilo2019identifyexoplanetsII}) focused on the challenge of vetting the large number of planet candidates identified in K2 light curves, which are especially noisy due to pointing instabilities from the Kepler telescope. This created difficulty in classifying true planet signals from false positives such as eclipsing binaries or systematics. To address this, they extended the AstroNet CNN first built for Kepler data by Shallue and Vanderburg (2018), adapting it to the shorter and noisier K2 campaigns. Their approach used both real K2 TCEs labeled as candidates, binaries, or junk, and tested on simulated ``Kepler-lite'' data to tune their model under similar conditions. Each TCE was represented by two views of the light curve, a global view showing the full orbital phase and a local view zoomed in on the transit. These were normalized and then passed through separate convolutional and max pooling layers before being merged into fully connected layers of 512 nodes, with additional scalar features such as planet-to-star radius ratio and impact parameter included to strengthen classification. The final output gave the probability of a candidate being a true planet. Their CNN achieved high accuracy averaging 98\% on test sets and successfully uncovered two new super-Earths within the K2 dataset. This work highlights the need for both global and local inputs including additional scalar features as it improves against noisy data while also improving the classification of planets and false positives.

Osborne et al. (\citeyear{Osborn2019rapidclassificationtesscandidates}) addressed the challenge of rapidly classifying the thousands of planet candidates generated monthly by TESS. To further the issue each month around 500 hours were required to select the best candidates from each month's findings. As any sane person would be able to tell, this is a bottleneck as it requires too many resources and people that could be used elsewhere. To resolve this issue however, Osborne and his colleagues trained a CNN model based on Shallue and Vanderburg (2018) and Ansdell et al. (2018) work, specifically the AstroNet and ExoNet CNN models. They trained and validated their CNN strictly on high-fidelity simulated TESS light curves (the TSOP-301 dataset), which provided known ground truths for both planets and false positives. Each candidate was represented by both a local and global view of its light curve, with centroid information also included to help spot background eclipsing binaries. Alongside this, during training, they included balanced batch sampling to counter class imbalance (only 14\% of candidates were true planets) and applied augmentation techniques to improve generalization and reduce overfitting. When tested, their model achieved strong results on their simulated set, with a 97\% average precision in binary classification and a 92\% accuracy in identifying planets in their multi-class classification. However, when using the real TESS set, the model's performance dropped, only being able to recover 61\% TOIs (TESS Objects of Interest) as planets. Though the model was still able to flag around 200 new TCEs as possible planets. Results highlight a reduction in human workload and planet discovery but also show a gap between simulated and real data performance in CNNs. This brings a need to build a model that is more generalizable capable of handling both simulated and observational dataset together, which we plan to explore in our project.

\subsubsection{CNN Classifier for Next Generation Transit Survey}
Chaushev et al. (\citeyear{Chaushev2019ngts}) focused on improving the speed and consistency of exoplanet detection in large transit surveys, particularly the Next Generation Transit Survey (NGTS), where manual vetting of light curves is reliable but extremely time-consuming. To automate this process, the authors introduced a convolutional neural network (CNN) called PlaNET, which aimed to reduce dependence on large labeled datasets by using a mix of real NGTS observations, unlabeled data, and simulated planetary transits. This combination of data helped the model learn from both real and synthetic signals while improving generalization to noisy and incomplete data. Each light curve was detrended, normalized, and phase-folded so that transits aligned, creating a consistent input for classification. Their architecture was inspired by AstroNet, featuring both global and local views of the light curve to capture the full orbital shape and transit depths. Additional stellar information, such as radius and orbital period, was included to help the model separate planets from false positives like eclipsing binaries or background blends. The model was trained using the Adam optimizer with early stopping and 10-fold cross-validation to avoid overfitting. In testing, PlaNET achieved around 96\% accuracy and an AUC of 99.3\%, successfully recovering 13 out of 14 confirmed NGTS planets while identifying several new candidates missed by human vetting. Their results show that CNNs can greatly reduce manual effort and improve the reliability of large-scale exoplanet searches. The main takeaway here is the use of extra features such as the radius and orbital period. These features further help in the process of classification, especially in the case where exoplanet lightcurves look similar to EBs and BEBs.

\subsubsection{Stella CNN Flare Classifier}
Feinstein et al. (\citeyear{Feinstein2020stella}) built a CNN called stella to detect stellar flares in TESS short-cadence data. The problem they faced was that manually identifying flares across thousands of young stars is too time-consuming, especially since high stellar variability makes simple outlier detection methods fail. Their solution was to train a CNN to recognize flare shapes directly from the light curves and use this to study how flare rates change with stellar age and temperature. They worked with 3,200 young stars from TESS Sectors 1–20, detrending and normalizing each light curve before splitting them into 200-cadence segments centered on either known flare events or non-flare regions. The architecture used standard components with one-dimensional convolutional layers to extract features, pooling layers to reduce dimensions, fully connected layers with dropout to prevent overfitting, and a final sigmoid layer for binary classification. They trained it with the Adam optimizer and binary cross-entropy loss, running 10 models with different initializations to create an ensemble and reduce overfitting. Their model stella performed really well, hitting an AUC of 0.993 and F1 scores above 0.98 on validation data. When applied to the full sample, it found over 17,000 flares and showed that flare activity drops off significantly for stars older than 50 Myr and hotter than 4000 K. What's interesting here is the data preparation approach. Instead of phase-folding like most exoplanet detection methods, they used fixed-length segments centered on events of interest. This could also work for planet detection, where we would just center segments on threshold crossing events instead of flares.

\subsubsection{GPU-accelerated Phase Folding}
Wang et al. (\citeyear{Wang2024gpfc}) developed GPFC, a pipeline that combines GPU-accelerated phase folding with a CNN to detect ultra-short-period planets with orbital periods less than one day. Most exoplanet detection pipelines rely on Box Least Squares (BLS) to identify transits, but BLS becomes extremely slow when scanning through thousands of trial periods at high resolution. Wang et al. took a different route by offloading the phase folding computation to the GPU and using a CNN to score each folded result. The workflow starts with detrending and outlier removal, then creates a dense grid of trial periods to test. Next, the light curve gets folded and binned at each period, followed by noise normalization to standardize the input. Finally, a CNN evaluates each folded light curve and outputs a score indicating the likelihood of a real transit signal. The impressive part is the speed at which they can fold a single light curve at 100,000 different trial periods in just five seconds on an RTX 3090 Ti. This works by batching 16-period folds simultaneously and using efficient binning algorithms that take advantage of GPU parallelization. The CNN itself has 19 layers and includes circular convolution, global max pooling, and dropout layers to prevent overfitting. They trained it using TensorFlow with the Adam optimizer. When they tested GPFC against traditional BLS methods, it performed better at detecting low signal-to-noise ratio transits and ran much faster. In a blind search over the 0.2-1.0 day period range, GPFC successfully recovered all confirmed ultra-short-period planets in the Kepler catalog. What makes this approach valuable is that GPU acceleration turns exhaustive period searches from impractical to practical. Testing 100,000 trial periods per light curve would take hours with the usual BLS method, but the GPU implementation cuts this down to seconds. Our current pipeline doesn't use GPU-accelerated folding yet, but this technique could become important if we need to process larger datasets or want to search finer period grids in future work.

\subsubsection{Applying These Ideas to Our Work}
Looking at what worked across these papers, we're going to borrow several techniques that consistently improve detection performance. The dual-view architecture from Shallue and Vanderburg (\citeyear{Shallue2018identifyexoplanetsI}) is one of the most useful ideas. It splits each light curve into a global view (2001 bins) covering the full orbital phase and a local view (201 bins) zoomed in on the transit itself. This way, the CNN gets both the orbital context and the finer transit details. We'll follow their preprocessing approach too: detrend with median subtraction, use Box Least Squares to find the period, phase fold around the transit center, and normalize so the baseline sits at zero with the deepest transit point at -1.

The auxiliary scalar features from Dattilo et al. (\citeyear{Dattilo2019identifyexoplanetsII}) and Chaushev et al. (\citeyear{Chaushev2019ngts}) seem worth including as well. Adding things like planet-to-star radius ratio, impact parameter, transit depth, and orbital period helps the model distinguish real planets from eclipsing binaries, especially when the light curves alone are ambiguous or noisy. Osborne et al. (\citeyear{Osborn2019rapidclassificationtesscandidates}) showed that centroid information can catch background eclipsing binaries that might otherwise slip through, and their balanced batch sampling approach handles the severe class imbalance problem we're likely to face. We should also test some of the data augmentation techniques that popped up across these studies to make the model more robust.

One issue that kept showing up in all these papers is the generalization problem. Models trained on real data often struggle with simulated datasets and vice versa. To address this, we're planning to train and test on both real observational data (Kepler, K2, TESS) and simulated light curves. That way, we can actually measure how well the model transfers across different missions and figure out what techniques improve cross-dataset performance.

\subsection{Recurrent Neural Networks (RNNs)}

This section reviews RNN applications to astronomical time series that inform the BiLSTM + clustering approach. We first summarize six foundational papers from our midterm report, then present three new studies that establish performance benchmarks and validate key architectural decisions.

\subsubsection{Foundational Work (Midterm References)}

\textbf{Vida et al.\ (2021)} trained stacked LSTMs to detect stellar flares in Kepler and TESS photometry, using dropout regularization and minority class weighting \cite{Vida2021findingflaresrecurrent}. \textit{Why I used it}: Their finding that models generalize well across Kepler and TESS supports consistent preprocessing, and their comparison showing stacked LSTMs outperform GRUs for astrophysical noise informed my choice of multi-layer BiLSTM architecture.

\textbf{K\"ugler et al.\ (2016)} paired Echo State Networks with autoencoders for unsupervised representation learning on Kepler light curves, evaluating reconstructions in sequence space rather than embedding space \cite{kugler2016explorativeapproachkepler}. \textit{Why I used it}: Their emphasis on sequence-level evaluation motivates evaluating my model on full windows and visualizing predictions across time to verify alignment with actual transit events.

\textbf{Du et al.\ (2016)} introduced Recurrent Marked Temporal Point Processes (RMTPP), demonstrating how RNNs can encode event timing history for predicting when and what type of event occurs next \cite{Du2016markedtemporalpoint}. \textit{Why I used it}: This concept motivates incorporating periodicity information---genuine transits repeat on predictable schedules while false positives are aperiodic---through BLS period estimates as clustering features.

\textbf{Speiser et al.\ (2020)} demonstrated that K-means clustering ($k$=5--20) followed by cluster-aware classification outperforms global classifiers by 8-15\% in F1 score on heterogeneous microscopy datasets \cite{Speiser2020clustering}. \textit{Why I used it}: This directly inspired my clustering approach---rather than training one BiLSTM to handle all stellar variability, I cluster windows by BLS features and let the network specialize via learned embeddings.

\textbf{Vu et al.\ (2024)} validated 3-layer stacked LSTM architectures for time series forecasting, showing 23\% improvement over 2-layer variants and 62\% F1 improvement from class weighting on imbalanced data \cite{Vu2024LSTM}. \textit{Why I used it}: Their findings directly inform my hyperparameter choices: 3-4 LSTM layers, dropout 0.3-0.4, and inverse class ratio weighting (pos\_weight = 3.367).

\textbf{Ding et al.\ (2024)} applied bidirectional LSTMs to photometric redshift estimation on 4.2 million SDSS galaxies, achieving 34\% outlier reduction versus MLP and demonstrating cross-survey transfer learning \cite{Ding2024photometric}. \textit{Why I used it}: Their success with BiLSTM for astronomical photometry and validation of transfer across surveys (SDSS to DES) supports my architecture choice and suggests Kepler-trained models can adapt to TESS.

\subsubsection{New Reference: Variable Star Classification with CNN-LSTM}

Becker et al.\ (2021) developed a hybrid CNN-LSTM architecture for classifying variable star light curves from the Optical Gravitational Lensing Experiment (OGLE) and Catalina Real-Time Transient Survey (CRTS), achieving 92-95\% accuracy across eclipsing binaries, RR Lyrae, Cepheids, and long-period variables \cite{Becker2021variablestars}. Their work directly addresses the challenge of classifying astronomical time series with irregular sampling and heterogeneous noise characteristics---precisely the challenges faced in exoplanet transit detection.

\textbf{Architecture and Design Rationale.} The hybrid combines 1D convolutional layers (kernel sizes 3, 5, and 7) with stacked LSTM layers in a sequential pipeline. The CNN component extracts local morphological features---transit shape, ingress/egress profiles, eclipse symmetry---through learned convolutional filters. Max pooling between convolutional blocks reduces sequence length from thousands of points to hundreds, improving computational efficiency while preserving salient features. The LSTM component then processes these compressed representations to capture long-range temporal dependencies including periodicity patterns, phase relationships between primary and secondary eclipses, and gradual brightness variations characteristic of pulsating stars.

The authors experimented with several architectural variants: CNN-only, LSTM-only, and the hybrid CNN-LSTM. The hybrid consistently outperformed single-architecture approaches by 5-8\% in classification accuracy, demonstrating that local feature extraction and sequential modeling provide complementary information. This finding directly informs my decision to use bidirectional LSTMs, which similarly combine forward and backward temporal context.

\textbf{Class Imbalance Handling.} Their dataset exhibited 10-20$\times$ class imbalance between common variable types (eclipsing binaries) and rare categories (anomalous Cepheids). They addressed this through stratified mini-batch sampling---ensuring each training batch contained representatives from all classes---combined with inverse-frequency class weights in the cross-entropy loss function. The weighted loss assigns higher penalties to misclassifying rare classes, preventing the model from simply predicting the majority class. This is precisely the strategy I implement with pos\_weight = 3.367 (later updated to 7.41 for Sector 1 data) to handle the 11.7\% positive class ratio in transit detection.

\textbf{Key Findings and Implications.} The hybrid achieved 92-95\% accuracy using raw photometric time series without hand-engineered period or amplitude features. This end-to-end learning approach validates my architectural choice to feed normalized flux windows directly into the BiLSTM rather than computing intermediate representations. The authors also demonstrated that their model generalized across OGLE and CRTS surveys despite different cadences, photometric systems, and noise characteristics---providing evidence that deep learning can learn mission-agnostic astronomical patterns.

\textbf{Relevance to My Implementation.} Two key lessons emerge: (1) raw time-series input works effectively, eliminating potential information loss from aggressive feature engineering; (2) weighted loss successfully handles severe class imbalance while maintaining high accuracy on minority classes. The hybrid architecture's success also motivates future work exploring CNN preprocessing before my BiLSTM, though the current architecture already achieves competitive results.

\subsubsection{New Reference: Machine Learning for Ground-Based Transit Surveys}

Schanche et al.\ (2019) developed an automated vetting pipeline for transit candidates from the Wide Angle Search for Planets (WASP) ground-based survey, achieving approximately 90\% correct identification of genuine planetary transits versus false positives \cite{Schanche2019mlexoplanet}. Ground-based surveys face substantially greater noise challenges than space-based missions due to atmospheric scintillation, weather interruptions, and day/night observing gaps, making their results particularly encouraging for space-based applications.

\textbf{The False Positive Problem.} A central challenge in transit detection is distinguishing genuine planetary transits from astrophysical false positives. Eclipsing binary stars produce transit-like dips but with several distinguishing characteristics: deeper eclipses (typically 1-50\% versus 0.01-1\% for planets), longer durations relative to orbital period, V-shaped rather than U-shaped profiles due to limb darkening effects, and often visible secondary eclipses when the smaller star passes behind the larger. Background eclipsing binaries---unrelated binary systems whose light contaminates the target aperture---are particularly insidious because their diluted signals can mimic planetary transit depths. Schanche et al.\ explicitly trained their classifiers to distinguish these scenarios.

\textbf{Methodology and Classifier Comparison.} The authors trained both Random Forest ensemble classifiers and convolutional neural networks on light curves containing four categories: confirmed planetary transits, eclipsing binaries (both foreground and background), intrinsic variable stars, and non-periodic systematic artifacts. Individual classifiers achieved 85-88\% accuracy but produced substantial false positive rates on challenging cases where transit-like signals had ambiguous characteristics. However, ensemble voting---combining predictions from multiple classifiers trained on different feature subsets---reduced errors by 15-20\%, achieving the final 90\% accuracy.

This ensemble finding resonates strongly with my clustering strategy. Rather than forcing a single model to handle all stellar variability types, K-means clustering partitions the data into homogeneous subgroups, and cluster embeddings allow the BiLSTM to specialize its learned representations for each group. The philosophical parallel is clear: both approaches recognize that a single classifier struggles with heterogeneous data, and stratification (whether through explicit ensembles or learned embeddings) improves performance.

\textbf{Discriminating Features.} The authors identified transit depth and duration as the two most important features for distinguishing planets from eclipsing binaries, with depth providing 3$\times$ more discriminative power than the next most important feature. Planets produce shallow transits (typically 0.01-1\% flux decrease) while binaries produce deeper eclipses. Transit duration relative to orbital period also differs: planetary transits last 1-5 hours while binary eclipses often span 5-15\% of the orbital period. These findings directly motivate my inclusion of BLS-derived depth and duration features in the clustering step, ensuring windows are grouped by these physically meaningful characteristics.

\textbf{Relevance and Performance Implications.} Ground-based surveys operate under conditions far noisier than space-based missions. Atmospheric turbulence introduces correlated noise, clouds create data gaps, and limited nightly observing windows reduce phase coverage. If machine learning achieves 90\% accuracy under these challenging conditions, my BiLSTM with cluster embeddings operating on cleaner TESS space photometry should achieve comparable or superior results. The emphasis on depth and duration as primary discriminators validates my BLS-feature-based clustering approach, confirming that these physically motivated features capture the essential differences between planet and non-planet signals.

\subsubsection{New Reference: Feature-Based Detection with Gradient Boosting}

Malik et al.\ (2022) took a fundamentally different approach to exoplanet detection: rather than feeding raw light curves into neural networks, they extracted 789 statistical features using the TSFRESH automated feature extraction library and trained a LightGBM gradient boosting classifier, achieving AUC 0.948 and recall 0.96 on Kepler data \cite{Malik2022exoplanetdetection}. This feature-based approach provides an important benchmark against which neural network methods must justify their added complexity.

\textbf{Feature Engineering with TSFRESH.} The Time Series Feature Extraction based on Scalable Hypothesis tests (TSFRESH) library automatically computes a comprehensive set of time series descriptors without requiring domain expertise. The 789 features span multiple categories: basic statistics (mean, variance, skewness, kurtosis), autocorrelation at various lags capturing periodicity, spectral properties from Fourier decomposition, entropy measures quantifying signal complexity, and change-point statistics identifying discontinuities. This exhaustive feature set describes the shape, periodicity, and statistical properties of each light curve without requiring the researcher to manually engineer representations.

The authors applied feature selection using a combination of statistical significance tests and recursive feature elimination, ultimately identifying approximately 150 features as most predictive. Top-ranked features included autocorrelation at lags corresponding to orbital periods, variance ratios between in-transit and out-of-transit segments, and spectral power at transit-related frequencies. Interestingly, many of these selected features correspond to the same physical information captured by BLS analysis---periodicity, transit depth, and duration---validating the domain-specific features I use for clustering.

\textbf{LightGBM Training and Results.} LightGBM, a gradient boosting framework optimized for speed and memory efficiency, constructs an ensemble of decision trees where each subsequent tree corrects errors made by previous trees. The authors trained on Kepler confirmed planets (positive examples) and a carefully curated set of non-planet signals including eclipsing binaries, variable stars, and systematic artifacts (negative examples). Five-fold cross-validation produced remarkably consistent results: AUC 0.948 $\pm$ 0.012, precision 0.91, recall 0.96.

The 0.96 recall indicates the classifier correctly identified 96\% of genuine planetary signals---missing only 4\% of planets in the dataset. The 0.948 AUC means that a randomly selected planet signal would be ranked above a randomly selected non-planet 94.8\% of the time. These results establish a concrete performance target that any neural network approach must meet or exceed to justify its additional architectural complexity.

\textbf{Benchmark Implications for Neural Networks.} Malik et al.'s results pose a critical question: if feature engineering followed by gradient boosting achieves AUC 0.948, what do neural networks offer? The answer lies in representation learning. Gradient boosting treats each of the 789 features independently, learning optimal decision boundaries in feature space but unable to capture complex interactions between sequential elements. My BiLSTM learns temporal patterns directly from raw flux---the characteristic U-shaped transit profile, the relationship between ingress and egress phases, subtle asymmetries indicating grazing transits---that feature-based methods cannot represent.

The architectural comparison reveals complementary strengths: feature-based methods excel when domain knowledge can be encoded into explicit features, while neural networks excel at discovering patterns that resist manual specification. My approach combines both philosophies: explicit BLS features (period, depth, duration) inform clustering, while the BiLSTM learns implicit temporal representations within each cluster. This hybrid strategy aims to capture the best of both worlds.

\textbf{Target Setting.} Malik et al.'s 0.948 AUC establishes the benchmark my BiLSTM must approach. Final results (AUC 0.9261 on TESS Sector 1 test data) approach this target, though several factors complicate direct comparison: different datasets (Kepler versus TESS), different class distributions (their balanced sampling versus my 11.1\% positive rate in test), and different evaluation protocols. The 0.9261 AUC represents a 33.3\% improvement over the baseline model (0.6947) and demonstrates competitive performance with feature-engineering approaches.

\subsubsection{Synthesis: Informing the BiLSTM + Clustering Design}

These nine works collectively inform the architecture and training strategy:

\begin{itemize}
    \item \textbf{Architecture}: 3-4 layer BiLSTM with 256 hidden units (Vida, Vu, Ding, Becker)
    \item \textbf{Clustering}: K-means ($k$=5) on BLS features before training (Speiser, Schanche)
    \item \textbf{Class balancing}: pos\_weight = 7.41 in BCEWithLogitsLoss (Vu, Becker)
    \item \textbf{Input}: Raw normalized flux windows (Becker)
    \item \textbf{Regularization}: Dropout 0.3-0.4 + early stopping (Vida, Vu)
    \item \textbf{Evaluation}: Sequence-level diagnostics (K\"ugler), timing features (Du)
    \item \textbf{Target}: AUC $>$ 0.948 (Malik)
\end{itemize}

This design addresses class imbalance through weighted loss, heterogeneous stellar variability through cluster embeddings, and temporal dependencies through bidirectional processing.

\subsection{Transformer Architecture}
Attention mechanisms have been found to perform very well in sequence modeling, allowing models to capture dependencies between tokens regardless of the distance between positions. Self-attention (or intra-attention) creates a relationship between different positions of the input by using an encoder-decoder structure made of stacked layers.

The research paper "Attention Is All You Need" \cite{vaswani2017attentionneed}, is the first of its kind to fully implement the Transformer architecture, and demonstrates this by training the model on English-to-German and English-to-French translation. Many of the researchers on this project were concurrently working at Google Brain and Google Research. This paper is currently one of the most influential of the 21st century and is listed as the 7th most-cited piece of research in the last 25 years by Nature. Transformer architecture has been the major basis for most Natural Language Processing (NLP) projects as well as generative Large Language Models (LLMs).

\citeauthor{vaswani2017attentionneed}'s design of Multi-Head Attention accomplishes this by projecting queries, keys, and values $h$ times in parallel. This allows the model to attend every token to all the other tokens in the sequence, and weights the positions with the average attention to other tokens, in each of the six layers of the encoder and decoder. Additionally, each head focuses on a different pattern/relationship in the data. The parallelization within the heads helped reduce the sequential limitations of RNNs when operating on larger data sets.

\citeauthor{vaswani2017attentionneed} made two models: a small one and a big one. When trying to test a translation from English to German, the big Transformer model achieved a new high BLEU score of 28.4, surpassing previous models by over 2.0 BLEU, and was trained on 8 NVIDIA P100 GPUs for 3.5 days. Even the small model, which was only trained for 12 hours, still beat the other models by a smaller margin. The large model also beat all other models when translating English to French, though only by 0.5 BLEU, and the smaller model preformed similarly to the others.

This experiment, though not focused on the interpretation of light curves, gives us an encouraging basis to pursue how effective an Attention-based model could be in analyzing transit data.

This idea was pursued by \citeauthor{morvan2022dontpayattentionnoise} in their \citeyear{morvan2022dontpayattentionnoise} paper, "Don't Pay Attention to the Noise: Learning Self-supervised Representations of Light Curves with a Denoising Time Series Transformer." In this paper, they designed a Transformer model named the Denoising Time Series Transformer (DTST) to try to remove noise, fill in missing data, and fine-tune the light curves to make it easier to analyze. The architecture uses real stellar light curves from TESS's first sector, showing greater adaptability than other traditional denoising techniques. \citeauthor{morvan2022dontpayattentionnoise} used the same positional embedding as \citeauthor{vaswani2017attentionneed} (\citeyear{vaswani2017attentionneed}) to keep track of the chronological sequence of points on the light curve.

After standardizing the light curve data from TESS and training the model, \citeauthor{morvan2022dontpayattentionnoise} found that results would not always denoise noisy data and might overfit quiet data, due to a lack of data normalization. The DTST research explicitly addressed the challenges of light curves, which are noted as "particularly challenging data objects." Although the results were not as consistent as desired, the experiment showed a lot of potential for the continued exploration of building a more robust Transformer model. They additionally found they could track how important the model thought different parts of the data were, by using an attention map.

The work "Distinguishing a planetary transit from false positives: a Transformer-based classification for planetary transit signals," \cite{salinas2023distinguishingtransitfalsepositives}, directly continues the search for efficient, interpretable models to classify TCEs and determine the possibility of whether the TCEs were actually exoplanets or a false positive. Other models typically used CNNs with many layers, which could quickly become computationally complex. \citeauthor{salinas2023distinguishingtransitfalsepositives}'s classification system was specifically designed to use phase-folded light curves from TESS, as well as flux data and centroid pixel position information, and included extra stellar and transit parameters (like period, duration, and host star temperature). The inputs were embedded into a vector representation using a convolutional layer before going to the multihead attention layer. Much of this preprocessing was handled with the Python packages `lightkurve` and `astropy api`. \citeauthor{salinas2023distinguishingtransitfalsepositives} ended up scaling the light curves "by removing the mean and scaling to unit variance for both centroids." When training the model, they only used one decoder which performed like an activation function, and three encoders.

Once trained, the Transformer achieved competitive results as other existing CNNs applied to TESS data with a similar goal, with average recall, precision, and F1 scores of 0.88. A major benefit of Transformers is the enhanced interpretation, and that the self-attention mechanism allowed the authors to observe that parameters like transit depth and host star temperature were highly influential in the classification decision.

Due to the continued results in pursuing Transformers, Salinas came back with another group of researchers to expand upon their previous research. This time the project was "Exoplanet transit candidate identification in TESS full-frame images via a transformer-based algorithm," \cite{salinas2025tessfullframe}. The research is focused around removing the need for highly pre-processed and phase-folded light curve data before analyzing for potential exoplanet candidates. \citeauthor{salinas2025tessfullframe} wanted to identify exoplanets that only had a single transit within the data segments (monotransiters), which would remove the periodic requirement from the light curve data. Because TESS data is limited to approximately 27 days, requiring a repeating transit will only find planets with orbits that take less than 25 days. This would additionally help identify plantets with Transit Timing Variations (TTVs) which could indicate a star system with multiple planets.

This approach continues the classification properties from the previous study \cite{salinas2023distinguishingtransitfalsepositives} and takes inputs of raw flux, centroid data, and background time-series from full-frame TESS images. Test data included labeled light curves that indicated confirmed exoplanets, confirmed eclipsing binaries or other false positives, and non-transit signals. The process begins by finding local features (short-range patterns) with a convolutional embedding layer (similar to \cite{salinas2023distinguishingtransitfalsepositives}), and these tokens are then fed into the multi-headed attention encoder to learn additional patterns.

Because the model can evaluate long-range dependencies, it can identify the characteristic shape of a planetary transit curve even if only a single transit is shows up in the entire 27 days of observation. This is important for detecting monotransiters and planets with longer periods.

When the model was tested on brand new TESS data, it found approximately 0.1 percent of the light curves indicated an exoplanet transit. With eight layers of encoding heads, they found an average Area Under Curve (AUC) score of 0.88 and an average F1 score of 0.83.

By the end of the experiment, \citeauthor{salinas2025tessfullframe} ended up testing sectors 1-26 of TESS data, which included 4.1 million light curves. Again, the model found approximately 0.1 percent to be potential candidates for planet transits. After confirming data through other means of verification, about 25\%, or 214, of the candidates were strong exoplanet candidates, including 88 single-transit candidates, which demonstrates the model's capability to detect transit events regardless of periodicity.

When looking at the paper "Transformer Architectures for Time Series Classification" \cite{wen2023transformerstimeseriessurvey}, we can see how the utilization of Transformer architectures in time series analysis stems from their fundamental ability to capture long-range dependencies and interactions in sequential data. This is particularly attractive for classification tasks, as Transformers have a prominent capability in capturing long-term dependency necessary for distinguishing sequence categories.

Recent research shows that optimizing the attention mechanism to capture temporal structure is key to superior performance. For instance, the MVGFormer model, which uses a visibility graph-guided attention mechanism (VG-Attention), excels in sequence-level classification \cite{chen2025graphbased}. In evaluations against existing methods on 17 multivariate time series datasets from UEA, MVGFormer achieved superior performance and demonstrated a strong representation capability in the vector space, effectively distinguishing between different categories. This success is attributed to its ability to focus attention on specific key time points or time intervals, such as periodic local modules or long-term trending structures, avoiding the dispersed attention characteristic of standard Transformer methods.

While standard Transformer models are adept at modeling global correlations, they often lack inherent assumptions about data locality. Conversely, traditional Convolutional Neural Networks (CNNs) are limited by their localized receptive fields, which makes them inefficient at capturing interactions between long-range information without stacking numerous layers. This limitation has often resulted in astronomic analyses focusing on the frequency domain (power spectra), which discards crucial non-Gaussian phase information present in the time domain.

To overcome these constraints and fully leverage the richness of long time series like stellar light curves, hybrid models that integrate both approaches are essential. The Astroconformer model provides a direct architectural precedent for my project. Astroconformer is a Transformer-based deep learning framework designed to capture both local short-range and global long-range information in stellar light curves, which contain information about oscillations and granulation \cite{pan2024astroconformer}.

The architecture successfully integrates the following components:

\begin{enumerate}
    \item {Convolutional Modules:} Astroconformer incorporates convolutional modules (including a convolutional layer, batch normalization, and Swish activation) to handle local feature extraction. This convolutional component is crucial because while the multi-head self-attention (MHSA) is theoretically sufficient, incorporating these convolutional layers enhances performance when labeled data is relatively limited.
    
    \item {Self-Attention Layers (MHSA):} The MHSA component is used to capture long-range dependencies.
\end{enumerate}

Ablation studies comparing Astroconformer variants confirm the strategic value of this hybrid approach. The data shows a clear performance decline as the receptive field size diminishes across various models (including pure CNNs like ResNet-18). Specifically, \citeauthor{pan2024astroconformer} found that replacing the MHSA module with a standard linear projection layer led to the most significant performance drop in forecasting tasks for MVGFormer, and in log $g$ inference for Astroconformer, underscoring the key role of the global attention mechanism in modeling long-range dependencies.

In the context of light curves, this architecture enables the model to leverage information beyond just the oscillations signal, successfully determining parameters like log $g$ based on granulation patterns, which is critical for dwarf and subgiant stars where oscillation signals may be brief or unresolved. The model achieved high precision, with an RMSE for log $g$ of 0.017 dex in dense sampling regimes.

The initial use of a CNN (or convolutional layer) for input preparation is also recommended. This embedding step aims to extract localized, short-range features before the Transformer analyzes the global sequence. Astroconformer similarly segments the input light curve into patches (e.g., 20 timestamps, representing 10 hours of data) and transforms them into feature vectors using an MLP (which acts as an embedding layer), effectively creating a coarse temporal axis before the encoder extracts local and global features. This is similar to \cite{salinas2023distinguishingtransitfalsepositives} where planetary transits in TESS data used a similar approach in which inputs were embedded into a vector representation using a convolutional layer before entering the multihead attention layer. 

\section*{Datasets}
For this research, we will be using both simulated and labeled observational data. The simulated data is available from the NASA Exoplanet Archive and includes injected on- and off-target transits, scrambled data, and false positives within K2 light curves. The labeled observational data will come from TESS datasets cross-referenced with matching labels, both from the Mikulski Archives for Space Telescopes (MAST). The purpose for choosing the above datasets is it provides a wide variety of light curve data taken from different sources. This means the light curves will have differing noise values and metrics to train from and test on. Furthermore, using these datasets will allow us to average our models' performances with a plethora of other research papers.

\begin{table}[t]
\centering
\begin{tabular}{p{0.28\linewidth} p{0.28\linewidth} p{0.33\linewidth}}
\hline
\textbf{Dataset} & \textbf{Purpose} & \textbf{Notes/SOTA} \\ \hline
Simulated Data ($\sim$10 known transits) & Train/val & Small, hand-checkable; class balance documented. Data from NASA Exoplanet Archive. \\
Pseudo/synthetic injections & Stress tests & Inject depth/duration/period jitter; report injection policy. Data from NASA \\
(Final Test) Kepler/TESS subsets & Generalization and best method determination & Add a SOTA line when datasets are finalized (TBD). Data from MAST and NASA Exoplanet Archive \\ \hline
\end{tabular}
\caption{{Planned datasets for ML learning and testing.}}
\label{{tab:datasets}}
\end{table}

\section*{Methodology}
\subsection*{Convolutional Neural Networks}
This work follows the dual-view CNN approach from Shallue and Vanderburg (\citeyear{Shallue2018identifyexoplanetsI}) and Osborne et al. (\citeyear{Osborn2019rapidclassificationtesscandidates}) for detecting exoplanets from light curves. The workflow has four main stages: preprocessing, model architecture, architecture/parameter tuning, and training.

\subsubsection{Preprocessing}
All light curves were processed using the AstroNet methodology. First, we cleaned the raw flux and time data with a $5\sigma$ clipping threshold to remove NaN values, infinities, and strong outliers. Then we applied a median filter with a kernel size of 401 to detrend long-term stellar variability while keeping the transit signal intact. We estimated the orbital period, mid-transit time ($t_0$), and duration using the Box Least Squares (BLS) algorithm, with Lomb–Scargle periodogram as a backup when BLS failed. Since BLS is computationally expensive, we had to use multi-threading to preprocess all 2821 light curves in a reasonable time. Without multi-threading, this would have taken around an hour or two, but with it, it is done within minutes, allowing for experimentation.

Next, each light curve was phase-folded around $t_0$ and normalized so the median out-of-transit baseline was zero and the minimum (transit depth) was –1. This follows what Shallue and Vanderburg (\citeyear{Shallue2018identifyexoplanetsI}) did and helps the CNN spot the transit curves better. We created two input representations: (1) a \textbf{global view} with 2001 bins covering the full orbital phase and (2) a \textbf{local view} with 201 bins centered on the transit. We used median binning instead of interpolation to keep the transit shape accurate and reduce outliers. Figure~\ref{fig:exoplanet_lightcurve} shows what a preprocessed light curve looks like after all preprocessing steps have concluded.

\begin{figure}[H]
  \centering
  \includegraphics[width=0.8\linewidth]{Images/CNN/cnn_preprocess_lightcurve.png}
  \caption{Preprocessed exoplanet light curve showing global and local views after phase folding and normalization.}
  \label{fig:exoplanet_lightcurve}
\end{figure}

\subsubsection{Model Architecture}
The classifier is a dual-branch 1D CNN that processes the global and local views separately before combining them at the end. Each branch has two convolutional blocks with GELU activations, Batch Normalization, and max pooling, followed by an adaptive pooling layer to get a fixed output size. The combined features go through three fully connected layers (512, 256, and 128 neurons) with dropout (0.5) before the final two-class softmax output. This setup lets the model learn both the orbital context from the global branch and the fine details of the transit from the local branch. Figure~\ref{fig:cnn_arch} shows the full layer structure.

\begin{figure}[H]
  \centering
  \includegraphics[width=0.7\linewidth]{Images/CNN/cnn_arch.png}
  \caption{Dual-branch CNN architecture showing how data flows from input through both branches to the final output.}
  \label{fig:cnn_arch}
\end{figure}

\subsubsection{Tuning}
Before training the full model, we needed to figure out the best hyperparameters for the job—learning rate, dropout rates, number of filters in each convolutional layer, etc. Manually testing all possible combinations would take forever, so we used Optuna, an automatic hyperparameter optimization framework. Optuna works by trying different parameter combinations, evaluating how well each one performs, and then suggesting new combinations to test based on what worked before.

We ran 50 optimization trials, where each trial trained a smaller version of the model for 10 epochs and reported the validation accuracy. The goal was to maximize accuracy while keeping training time down to a minimum. Optuna tested different architectures (number of convolutional layers, filter sizes, fully connected layer dimensions) and training parameters (learning rate, weight decay, dropout rate, batch size). Table~\ref{tab:architecture_params} shows the architecture parameters that Optuna searched through, and Table~\ref{tab:training_params} lists the training hyperparameters.

The tuning process helped us avoid a lot of guesswork and saved us a lot of time. Without it, we would've had to manually try dozens of configurations and probably would've missed the optimal setup. Our initial manual attempts only hit about 54\% accuracy, but the final hyperparameters from Optuna—learning rate of $1\times10^{-3}$, weight decay of $1\times10^{-4}$, dropout of 0.5, and batch size of 32—improved validation performance.

\begin{table}[H]
\centering
\caption{Architecture parameters optimized by Optuna}
\label{tab:architecture_params}
\begin{tabular}{lll}
\hline
\textbf{Parameter} & \textbf{Search Range} & \textbf{Optimal Value} \\
\hline
Number of Conv Layers & 1--3 & 2 \\
Conv1 Filters (Global) & 16--128 & 64 \\
Conv2 Filters (Global) & 32--256 & 128 \\
Conv1 Filters (Local) & 16--128 & 32 \\
Conv2 Filters (Local) & 32--256 & 64 \\
Kernel Size & 3, 5, 7 & 5 \\
FC Layer 1 Size & 256--1024 & 512 \\
FC Layer 2 Size & 128--512 & 256 \\
FC Layer 3 Size & 64--256 & 128 \\
Dropout Rate & 0.3--0.7 & 0.5 \\
\hline
\end{tabular}
\end{table}

\begin{table}[H]
\centering
\caption{Training hyperparameters optimized by Optuna}
\label{tab:training_params}
\begin{tabular}{lll}
\hline
\textbf{Parameter} & \textbf{Search Range} & \textbf{Optimal Value} \\
\hline
Learning Rate & $1\times10^{-5}$ to $1\times10^{-2}$ & $1\times10^{-3}$ \\
Weight Decay & $1\times10^{-6}$ to $1\times10^{-3}$ & $1\times10^{-4}$ \\
Batch Size & 16, 32, 64, 128 & 32 \\
Optimizer & Adam, AdamW, SGD & AdamW \\
\hline
\end{tabular}
\end{table}

\subsubsection{Training and Validation}
We trained the model using AdamW optimizer with a learning rate of $1\times10^{-3}$ and weight decay of $1\times10^{-4}$—these values came from Optuna hyperparameter tuning. The loss function was cross-entropy with class weights to handle the imbalance between planets and false positives. 

To help the model generalize better, we applied random augmentations to each light curve during training. These included Gaussian noise, random flux scaling, baseline offsets, phase shifts, variable transit depths, and partial bin dropout. The goal was to simulate realistic variations and missing data points that you see in actual TESS and Kepler data. Every time a light curve gets selected from a training batch, a different random augmentation is applied, so the same light curve looks slightly different each time the model sees it. This forces the model to learn more defined features instead of memorizing specific examples.

We used a learning rate scheduler (ReduceLROnPlateau) to reduce the learning rate when validation loss stopped improving, and gradient clipping ($\max norm=1.0$) to prevent training instability. The model was trained at a maximum of 60 epochs and only trained up to 18 epochs, with early stopping after 15 epochs of no improvement (Figure~\ref{fig:train-loss}). We tracked accuracy, precision, recall, F1-score, and AUC during validation and saved the best-performing model for testing.

\subsection{Recurrent Neural Networks}
\subsubsection{Problem Formulation}
How do we represent exoplanet transit detection as a machine learning problem? The task is binary classification of time-series windows. Given a 2048-point segment of a stellar light curve (normalized flux measurements over time), predict whether it contains a planetary transit (positive class) or represents non-planetary variability such as flares, noise, or eclipsing binaries (negative class). The input is a univariate time series $\mathbf{x} = (x_1, x_2, \ldots, x_{2048})$ where each $x_i$ represents normalized flux at time $i$. The output is a probability $p \in [0,1]$ indicating the likelihood that the window contains a transit. This formulation allows us to apply standard deep learning classification techniques to the astronomical time series problem.

\subsubsection{Data Collection and Preprocessing}
\textbf{Raw Data Sources.} The primary dataset is the TESS Sector 1 Ground Truth dataset, containing 13,541 labeled light curves across four categories: confirmed/candidate exoplanet hosts (3,146), non-planet stellar sources (8,624), eclipsing binaries (900), and background eclipsing binaries (871). This ground truth dataset was curated by the TESS team and provides reliable labels for supervised learning. Using real labeled data from a single TESS sector ensures consistency in noise characteristics and observation cadence while providing sufficient diversity in stellar types and transit signatures.

\textbf{Preprocessing Pipeline.} Each raw light curve undergoes several preprocessing steps to prepare it for the neural network. First, I remove obvious outliers using sigma clipping (points beyond 5$\sigma$ from the median). This eliminates spurious measurements that could confuse the model. Second, I apply a simple detrending procedure to remove long-term stellar variability using a median filter with a 101-point window. Why a median filter? This preserves the short-term transit dips while removing trends from stellar rotation or instrument drift that are not relevant to transit detection. Third, I normalize each light curve to zero mean and unit variance. This step is critical for training because TESS and Kepler have different noise characteristics, and normalization allows the model to learn features that generalize across missions rather than overfitting to mission-specific noise patterns. Fourth, I handle gaps in the data by simple linear interpolation for small gaps (fewer than 10 points) or flagging long gaps for exclusion during window extraction. This prevents the model from learning to recognize data gaps as transit features.

\subsubsection{Window Building and Feature Extraction}
\textbf{Sliding Window Approach.} I extract 2048-point windows from each preprocessed light curve, generating 3 windows per light curve with random starting positions. This produces 33,051 total windows from 13,541 light curves, split 80/20 into 26,472 training and 6,579 test windows. Why 2048 points? This window length corresponds to approximately 40-50 hours of observations at typical TESS cadence (2 minutes), sufficient to capture 2-3 transits for short-period planets. According to Kepler's third law, most detectable exoplanets have periods ranging from 1-50 days, so this window size balances temporal coverage with computational efficiency.

\textbf{Statistical Features for Clustering.} To prevent data leakage (using label-correlated features like transit depth), I compute eight statistical features from each window's flux values: mean, standard deviation, variance, skewness, range (max - min), median, median absolute deviation (MAD), and peak-to-peak amplitude. These features characterize the window's statistical properties without explicitly detecting transits, ensuring the clustering step remains unsupervised with respect to the transit labels.

\textbf{K-means Clustering.} I apply K-means clustering ($k=5$) to the eight-dimensional statistical feature space to group windows with similar noise and variability characteristics. Why clustering? This provides two benefits that directly improve model performance. First, it reduces effective problem complexity by allowing the model to learn cluster-specific patterns. Windows from quiet stars with low variability cluster separately from active stars with flares or strong rotation signals. Rather than forcing the BiLSTM to learn one global representation, clustering enables specialization. Second, it provides interpretability by grouping similar stellar types together. The cluster assignments are embedded into a 32-dimensional space and fed to the neural network, allowing it to learn relationships between clusters.

\subsubsection{Neural Network Architecture}
\textbf{BiLSTM Layers.} The core of my architecture is a 4-layer Bidirectional Long Short-Term Memory (BiLSTM) network. Each layer has 256 hidden units in each direction (512 total per layer). Why bidirectional? BiLSTM processes the input sequence in both forward and backward directions, allowing the model to use context from both past and future time points. This is important for transit detection because the shape of the ingress and egress (beginning and end of the transit) provides critical information about the planet-to-star radius ratio and orbital inclination. Bidirectional processing captures this better than forward-only LSTM, which only sees the ingress when it predicts the transit center.

Why 4 layers? According to deep learning theory, stacked LSTM layers learn hierarchical temporal representations at different scales. The first layer learns low-level features like local flux variations and noise patterns. The second layer combines these into mid-level features like the characteristic U-shaped transit pattern. The third and fourth layers capture high-level temporal structure like period detection across multiple transits within a window. Hyperparameter optimization using Optuna determined that 4 layers with 256 hidden units provided the best balance between capacity and generalization on the 26,472-window training set. The larger dataset (compared to initial experiments with 655 windows) supports deeper architectures without overfitting.

\textbf{Cluster Embedding Layer.} I embed the cluster assignments into a 32-dimensional continuous space using an embedding layer. Why embeddings instead of one-hot encoding? This allows the model to learn relationships between clusters rather than treating them as independent categories. For example, clusters containing short-period and medium-period planets might share some features (periodicity, U-shape), so the embedding space can represent this similarity through vector proximity. The embedding vector is concatenated with the final BiLSTM output (512 dimensions from forward and backward directions) to form a 544-dimensional representation that combines temporal features from the LSTM with cluster context from the embeddings.

\textbf{Fully Connected Classification Head.} The combined representation passes through three fully connected layers with batch normalization and ReLU activation. The first layer projects from 544 to 256 dimensions, compressing the information. The second layer projects from 256 to 128, further compressing. The final layer projects from 128 to 1 with sigmoid activation, producing a probability. Why batch normalization? It stabilizes training by normalizing layer inputs, allowing higher learning rates and faster convergence. Dropout with $p=0.311$ (optimized via Optuna) is applied after each hidden layer to prevent overfitting. The network has approximately 3.9 million trainable parameters, appropriate for the 26,472-window training set.

\subsubsection{Training Procedure}

\textbf{Loss Function and Class Balancing.} I use binary cross-entropy with logits loss (BCEWithLogitsLoss in PyTorch). To handle the 11.9 percent positive class imbalance (3,147 planets vs 23,325 non-planets in the training set), I apply a positive weight of 7.41 in the loss function. What does this mean? Errors on positive examples are weighted approximately 7.4 times higher than errors on negative examples, which encourages the model to prioritize recall (finding planets) over precision (avoiding false positives). This weighting was critical for preventing the model from collapsing to the trivial solution of predicting "no planet" for all inputs, which would achieve 88 percent accuracy but miss every real planet.

\textbf{Optimizer and Learning Rate Schedule.} I use the AdamW optimizer with a learning rate of $1 \times 10^{-4}$ and weight decay of $1 \times 10^{-5}$. Why AdamW rather than standard Adam? AdamW decouples weight decay from the gradient-based optimization, which improves generalization by preventing the optimizer from overfitting to training noise. I apply cosine annealing to the learning rate over 60 epochs, smoothly reducing it from the initial value to near zero. This schedule allows aggressive learning early in training when gradients are large and fine-tuning later when the model approaches a local minimum.

\textbf{Regularization.} Several techniques prevent overfitting. First, dropout ($p=0.311$, optimized via Optuna) randomly disables approximately 31 percent of neurons during training, forcing the network to learn redundant representations rather than memorizing specific training examples. Second, gradient clipping at norm 1.0 prevents exploding gradients in the recurrent layers, which is a common problem with deep LSTMs. Third, early stopping with patience 15 terminates training if validation AUC does not improve for 15 consecutive epochs, preventing the model from overfitting after it has learned the generalizable patterns. Fourth, I use mixed precision training (FP16) which provides a 1.6$\times$ speedup with minimal accuracy loss compared to full precision (FP32). Batch size was empirically determined to be 64 through benchmarking; larger batch sizes (128+) caused GPU memory pressure and slower training despite fewer iterations per epoch.

\textbf{Data Split.} The TESS Sector 1 ground truth dataset contains 33,051 windows extracted from 13,541 light curves. I use an 80/20 split: 26,472 windows for training and 6,579 windows for testing. The split is stratified by label to maintain the 11.9 percent positive ratio in both sets, ensuring the test set is representative. Unlike my early experiments with 655 windows from a curated sample, this dataset provides unbiased coverage of all stellar types observed in Sector 1, including challenging cases like eclipsing binaries and background contamination that could produce false positive transit signals.

\subsection{Transformer Architecture}
\subsubsection*{Data Description and Preprocessing}
The original dataset for the Transformer architecture began with 100 preprocessed light curves. These were divided into a set of 70 training samples, 15 validation samples, and 15 test samples. However, the dataset is in process of being increased to 3000 TESS observations and consists of 2,100 training samples, 450 validation samples, and 450 test samples, each comprising sequential flux measurements representing stellar brightness over time. Each light curve segment is padded or truncated to a fixed length of 4,320 observations (equivalent to 90 days of 30-minute cadence data), ensuring uniform input dimensions for the model.

Preprocessing on the data includes:

\begin{itemize}
    \item \textbf{Outlier Removal:} Sigma-clipping with a threshold of 5.0 standard deviations applied using rolling windows to mitigate spurious flux readings.
    \item \textbf{Trend Removal:} Low-frequency stellar variability trends are extracted via Savitzky-Golay filtering and removed by division of the flux, preserving subtle transit depths.
    \item \textbf{Interpolation:} Short temporal gaps smaller than 2 days are linearly interpolated to maintain temporal continuity.
    \item \textbf{Normalization:} Flux values are scaled to the \([0,1]\) range using min-max normalization, producing mean flux near 0.48 with a standard deviation around 0.20, facilitating stable and efficient optimization.
    \item \textbf{Segmentation:} Longer light curves are partitioned into fixed-length, optionally overlapping segments for batch processing during training.
\end{itemize}

Data augmentation heuristics such as noise addition and time warping were employed on the initial small dataset to improve model generalization against observational noise. The augmentation will be added to the larger dataset once optimized hyperparameters have been chosen.

\subsubsection*{Model Architecture}

The classification transformer model includes:

\begin{itemize}
    \item A \textbf{convolutional embedding} block composed of two 1D convolutional layers with ReLU activations and batch normalization to capture local temporal features.
    \item A \textbf{temporal positional encoding} submodule that utilizes actual observation timestamps, normalized to \([0,1]\), projected through a linear layer to impart sequential context to the transformer.
    \item A stack of \(n_{layers}=4\) transformer encoder blocks, each implementing multi-head self-attention with \(n_{heads}=4\) and feed-forward networks of dimensionality \(d_{ff}=512\) to model global temporal dependencies.
    \item Adaptive average pooling to condense sequence embeddings into fixed-dimensional representations.
    \item A classification head consisting of fully connected layers with non-linearities and dropout, outputting logits across two classes (transit/non-transit).
\end{itemize}

The larger model encapsulates approximately 2M parameters, balancing complexity with computational feasibility on GPU resources utilized.

\subsubsection*{Training Procedure}

The training regimen incorporates:

\begin{itemize}
    \item Weighted cross-entropy loss augmented with label smoothing (e.g., smoothing factor \(\approx 0.03\)) to reduce model overconfidence and improve calibration.
    \item The AdamW optimizer with learning rates explored in the range \([1\times10^{-5}, 1\times10^{-3}]\), employing weight decay around \(3.8 \times 10^{-4}\) for regularization.
    \item A warmup scheduler effectuating a linear increase in learning rate over 5 epochs, transitioning to cosine annealing for gradual decay.
    \item Gradient clipping using norms around 0.8 to stabilize backpropagation.
    \item Early stopping criteria with patience of roughly 10 epochs to preclude overfitting.
\end{itemize}

Batch sizes are kept moderate (e.g., 16) due to memory constraints on the allocated GPU environment (Google Colab).

\subsubsection*{Hyperparameter Optimization}
Hyperparameters are optimized via Optuna, leveraging Bayesian optimization via the Tree-structured Parzen Estimator (TPE) sampler. Trials are pruned based on intermediate validation loss to accelerate convergence toward optimal configurations. Key hyperparameters tuned include model depth, embedding dimensions, dropout rates, learning rates, and batch size.

\section*{Results}
\subsection {CNN + Augmentation Approach}
The trained CNN was tested on a preprocessed dataset of 2821 labeled samples (TESS Lilith 4 Dataset) split into 1974 for training, 423 for validation, and 424 for testing. The final model used two convolutional layers with 24 filters, kernel size of 5, Batch Normalization, and three fully connected layers with dropout.

\subsubsection{Training Performance}
Training ran for up to 60 epochs with class-weighted loss to handle the imbalance between planets (52.4\%) and non-planets (47.6\%). The model converged early and stopped at epoch 18 when validation loss stopped improving.

Figure~\ref{fig:train-loss} shows how training and validation loss changed over time, while Figure~\ref{fig:val-metrics} shows F1 and AUC performance. Validation accuracy settled around 74\%, with F1 around 0.79 and AUC between 0.77 and 0.79. These were the best results we obtained during training.

\begin{figure}[ht]
  \centering
  \includegraphics[width=0.8\linewidth]{Images/CNN/train_loss_curves.png}
  \caption{Training and validation loss over epochs. The model converged by epoch 18 with a validation loss of around 0.53.}
  \label{fig:train-loss}
\end{figure}

\begin{figure}[ht]
  \centering
  \includegraphics[width=0.8\linewidth]{Images/CNN/f1_auc_curve.png}
  \caption{Validation F1 and AUC scores over epochs. Performance stayed between 0.78 and 0.80 before early stopping.}
  \label{fig:val-metrics}
\end{figure}

The validation metrics stayed consistent across epochs with only small variations, and early stopping prevented overfitting. One good sign is that training and validation accuracy were really close at each epoch, having only a 3 to 7 percent difference. In earlier versions of this model, we saw gaps of up to 40\% by the end, which was a huge red flag for overfitting. Getting that down shows that the regularization is working as intended.

\subsubsection{Test Set Evaluation}
After the training finished, we tested the best model on the test set. Table~\ref{tab:test-metrics} has the final numbers.

\begin{table}[ht]
\centering
\small
\caption{Final test set performance of the CNN classifier.}
\label{tab:test-metrics}
\begin{tabular}{lcc}
\hline
\textbf{Metric} & \textbf{Score} & \textbf{Interpretation} \\
\hline
Accuracy   & 72.64\% & Overall correctness across both classes \\
Precision  & 0.647 & Fraction of predicted planets that were correct \\
Recall     & 0.864 & Fraction of true planets correctly identified \\
F1 Score   & 0.740 & Mean of precision and recall \\
AUC        & 0.816 & Ability to rank planets above false positives \\
\hline
\end{tabular}
\end{table}

\subsubsection{Analysis of Results}

The model got 86.4\% recall, which means that it has a low false-negative rate and catches most planetary transits. The precision was lower at 64.7\%. While precision is low, it's better to flag more candidates (including some false positives) than to miss real planets. We'd rather the model say "yes, possible planet" and be wrong sometimes than say "no" and throw away a possible discovery. Figure~\ref{fig:confusion_matrix} shows this. We have high counts for both true positives and true negatives, but also some false positives where the model thought a non-planet was a planet. This is a good thing, as missing a real planet is way worse than flagging a false positive that gets filtered out later. The AUC of 0.816 shows the model is decently good at ranking real planets higher than false positives. The F1 score of 0.740 shows a reasonable balance between precision and recall, which matches what other CNN approaches like AstroNet and the TESS versions got.

\begin{figure}[ht]
  \centering
  \includegraphics[width=0.8\linewidth]{Images/CNN/confusion_matrix.png}
  \caption{Confusion matrix showing true/false positive and negative rates.}
  \label{fig:confusion_matrix}
\end{figure}

Overall, the model learned to recognize transit features from the dual-view inputs. These results show that the AstroNet-style preprocessing and dual-branch architecture work well on unseen data while keeping recall high, which makes it useful for automated exoplanet candidate screening. Now, while the current model is at a good baseline, it can be improved in many ways to increase overall results and reliability.

\subsection{BiLSTM + Clustering Approach}

I trained a 4-layer BiLSTM with K-means clustering ($k$=7) on the TESS Sector 1 ground truth dataset. Table~\ref{tab:bilstm_dataset} summarizes the data: 33,051 windows from 13,541 light curves with 11.9\% positive class ratio. The model (3.07M parameters) uses 192 hidden units per direction, dropout 0.334, and positive class weighting (7.41) to handle imbalance. Hyperparameters were optimized via Optuna over 30 trials.

\begin{table}[h]
\centering
\caption{BiLSTM Dataset and Configuration}
\label{tab:bilstm_dataset}
\begin{tabular}{@{}lr|lr@{}}
\toprule
\textbf{Dataset} & \textbf{Value} & \textbf{Architecture} & \textbf{Value} \\
\midrule
Total Windows & 33,051 & BiLSTM Layers & 4 \\
Positive Class & 11.9\% & Hidden Units & 192 \\
Training Set & 26,472 & Parameters & 3.07M \\
Test Set & 6,579 & Dropout & 0.334 \\
Window Length & 2048 & Pos Weight & 7.41 \\
\bottomrule
\end{tabular}
\end{table}

\subsubsection{Final Test Results}

The final model achieved strong performance on the held-out test set of 6,579 windows. Table~\ref{tab:final_test_results} shows the key metrics, and Table~\ref{tab:confusion_matrix_rnn} presents the confusion matrix.

\begin{table}[h]
    \centering
    \caption{Final Test Set Performance (6,579 windows)}
    \label{tab:final_test_results}
    \begin{tabular}{@{}lr@{}}
        \toprule
        \textbf{Metric} & \textbf{Value} \\
        \midrule
        AUC & \textbf{0.9261} \\
        F1 Score & 0.5708 \\
        Precision & 39.93\% \\
        Recall & \textbf{100\%} \\
        Accuracy & 83.26\% \\
        \bottomrule
    \end{tabular}
\end{table}

\begin{table}[h]
    \centering
    \caption{Confusion Matrix - BiLSTM Test Set}
    \label{tab:confusion_matrix_rnn}
    \begin{tabular}{@{}lrr@{}}
        \toprule
         & \textbf{Pred. Positive} & \textbf{Pred. Negative} \\
        \midrule
        \textbf{Actual Positive} & TP = 732 & FN = 0 \\
        \textbf{Actual Negative} & FP = 1,101 & TN = 4,746 \\
        \bottomrule
    \end{tabular}
\end{table}

The 100\% recall is critical for exoplanet detection: the model detects every planet in the test set with zero false negatives. The lower precision (39.93\%) indicates some false positives, but missing a real planet is far worse than flagging a candidate for follow-up verification. The AUC of 0.9261 represents a \textbf{33.3\% improvement} over the baseline model (0.6947).

\subsubsection{Progressive Improvement History}

Table~\ref{tab:optimization_results} shows the progression from baseline to final model across the project timeline.

\begin{table}[h]
    \centering
    \caption{Model Performance Progression}
    \label{tab:optimization_results}
    \begin{tabular}{@{}lrrr@{}}
        \toprule
        \textbf{Version} & \textbf{Date} & \textbf{AUC} & \textbf{Notes} \\
        \midrule
        Baseline (655 windows) & Oct 2025 & 0.6947 & Initial model \\
        Optuna optimized & Nov 2025 & 0.7572 & +9.0\% \\
        Sector 1 (33K windows) & Dec 2025 & \textbf{0.9261} & \textbf{+33.3\%} \\
        \bottomrule
    \end{tabular}
\end{table}

\subsubsection{SMOTE Class Balancing Experiment}

I compared three class imbalance strategies: loss weighting (pos\_weight=3.367), naive up-sampling by duplication, and SMOTE interpolation. Table~\ref{tab:smote_comparison} and Figure~\ref{fig:smote_comparison} show SMOTE achieved the best results: AUC 0.8175 (+8.0\% over baseline), F1 0.7209, and 19/300 real planets detected. Naive up-sampling catastrophically failed (0/300 detections) despite good training metrics---it memorized duplicates rather than learning generalizable patterns. SMOTE's k-nearest neighbor interpolation creates diverse synthetic examples that preserve transit characteristics while introducing controlled variability.

    \begin{table}[h]
        \centering
        \small
        \caption{Class Imbalance Strategy Comparison}
        \label{tab:smote_comparison}
        \begin{tabular}{@{}lrrrrr@{}}
            \toprule
            \textbf{Strategy} & \textbf{Size} & \textbf{AUC} & \textbf{F1} & \textbf{Detected} & \textbf{Status} \\ \midrule Class Weight. & 655 & 0.7572 & 0.5145 & 16/300 & Good \\ Naive Up-samp. & 600 & 0.7210 & 0.6136 & 0/300 & Failed \\
            \textbf{SMOTE} & 600 & \textbf{0.8175} & \textbf{0.7209} & \textbf{19/300} & \textbf{Best} \\
            \bottomrule
        \end{tabular}
    \end{table}

    \begin{figure}[h]
        \centering
        \includegraphics[width=\columnwidth]{Images/RNN/bilstm_class_imbalance_comparison.png}
        \caption{Class imbalance strategy comparison. SMOTE achieved highest AUC (0.818) and detected 19/300 real planets (6.3\%) while naive up-sampling failed completely (0/300).}
        \label{fig:smote_comparison}
    \end{figure}
\subsubsection{Cross-Mission Generalization Analysis}

To test whether the TESS-trained model generalizes to other missions, I evaluated it on Kepler exoplanet hosts. The key challenge: TESS uses 2-minute cadence while Kepler used 30-minute cadence (15$\times$ difference), fundamentally changing transit signal morphology.

\textbf{Initial Evaluation.} On 279 windows from 93 confirmed Kepler planets, the TESS-trained model predicted all windows as positive (100\% recall, mean probability 0.978). However, without negative examples, precision and AUC could not be computed.

\textbf{Fine-Tuning Experiment.} Following professor recommendation, I tested whether fine-tuning on mixed TESS+Kepler data would improve cross-mission generalization. I downloaded 398 additional Kepler light curves from MAST, built 1,473 training windows from 491 total Kepler sources, and fine-tuned the model with 95\% TESS + 5\% Kepler data (learning rate $10^{-5}$, 20 epochs).

\begin{table}[h]
\centering
\caption{Fine-Tuning Results on Kepler Planet Detection}
\label{tab:finetuning_results}
\begin{tabular}{@{}lrrr@{}}
\toprule
\textbf{Model} & \textbf{Kepler Mix} & \textbf{Detection Rate} & \textbf{Mean Prob} \\
\midrule
Original TESS & 0\% & 55/1473 (3.7\%) & 0.156 \\
Fine-tuned & 5\% & 0/1473 (0\%) & 0.071 \\
\bottomrule
\end{tabular}
\end{table}

\textbf{Key Finding: Fine-tuning made performance worse, not better.} The original TESS model detected 3.7\% of Kepler planets; after fine-tuning with 5\% Kepler data, detection dropped to 0\%. Adding Kepler data caused the model to become more conservative overall, reducing mean probability across all test samples. The likely mechanism: the model learned that Kepler's 30-minute cadence produces smoother, less sharply-defined transit shapes than TESS's 2-minute cadence. When exposed to both during fine-tuning, the model could not reconcile these conflicting representations and instead learned to suppress predictions on data that did not match TESS characteristics precisely.

\textbf{Conclusion.} Simple fine-tuning does not solve the cross-mission generalization problem. The BiLSTM+Clustering architecture learns mission-specific features (cadence patterns, noise characteristics) rather than generalizable transit physics. This is a fundamental limitation requiring architectural changes such as cadence-agnostic preprocessing, domain-adversarial training, or multi-task learning with mission identification.

\subsection{Transformer Approach}
\subsubsection*{Model Configuration}
The evaluated model employs a  transformer architecture for use on a moderate dataset (2100 light curves) with the following specifications:
\begin{itemize}
    \item Embedding dimension: \(d_{model} = 64\)
    \item Number of encoder layers: \(n_{layers} = 2\)
    \item Number of attention heads: \(n_{heads} = 4\)
    \item Best validation loss: \(0.6263\)
\end{itemize}

\subsubsection*{Performance Metrics}

The model was evaluated on a held-out test set of 450 samples, achieving the following metrics:

\begin{table}[h]
\centering
\caption{Overall Classification Metrics}
\begin{tabular}{@{}lr@{}}
\toprule
\textbf{Metric} & \textbf{Score} \\
\midrule
Accuracy        & 0.6333 \\
Precision       & 0.6437 \\
Recall          & 0.6737 \\
F1-Score        & 0.6584 \\
ROC-AUC         & 0.6766 \\
\bottomrule
\end{tabular}
\end{table}

\subsubsection*{Per-Class Performance}

The classification report reveals relative symmetry between the two classes:

\begin{table}[h]
\centering
\caption{Classification Report by Class}
\begin{tabular}{@{}lcccc@{}}
\toprule
\textbf{Class} & \textbf{Precision} & \textbf{Recall} & \textbf{F1-Score} & \textbf{Support} \\
\midrule
Non-Transit & 0.6207 & 0.5888 & 0.6043 & 214 \\
Transit     & 0.6437 & 0.6737 & 0.6584 & 236 \\
\midrule
Macro avg.  & 0.6322 & 0.6313 & 0.6314 & 450 \\
Weighted avg. & 0.6328 & 0.6333 & 0.6327 & 450 \\
\bottomrule
\end{tabular}
\end{table}


\begin{figure}
    \centering
    \includegraphics[width=1.0\linewidth]{Images/RNN/confusion_matrix.png}
    \caption{Confusion matrix for most recent trial of Transformer training}
    \label{fig:placeholder}
\end{figure}

\subsubsection*{Analysis}

The model shows stronger performance on the Transit class with precision of 0.6437, recall of 0.6737, and F1-score of 0.6584. The moderate recall (67.37\%) indicates the model successfully identifies many actual transit events. The precision shows that when the model predicts a transit, it's correct about 64\% of the time.

Performance is slightly weaker for Non-Transit events with precision of 0.6207, recall of 0.5888, and F1-score of 0.6043. The lower recall (58.88\%) means the model misclassifies more non-transit light curves as transits, contributing to false positives.

The macro average (0.6314) and weighted average (0.6327) F1-scores are nearly identical, confirming that both classes contribute equally to overall performance. The model doesn't show strong bias toward either class.

The 63\% performance across all metrics suggests the model has learned meaningful patterns in the light curve data but still struggles with approximately 37\% of cases. This indicates more room for improvement in the model. 

\subsubsection*{Future Improvements}

To improve model performance:
\begin{enumerate}
    \item Increase model capacity by expanding \(d_{model}\) and \(n_{layers}\).
    \item Extend hyperparameter tuning to explore larger learning rates and regularization parameters.
    \item Investigate more ways to help the model train on the data.
\end{enumerate}


\section*{Conclusion}

This work investigated three deep learning architectures---CNN, RNN (BiLSTM), and Transformer---for automated exoplanet transit detection in TESS light curve data. Each team member developed and optimized a distinct architecture, enabling direct comparison of their strengths and limitations on the same TESS Sector 1 ground truth dataset.

\textbf{RNN Findings.} The BiLSTM architecture with K-means clustering achieved \textbf{AUC 0.9261} on the 6,579-window test set with \textbf{100\% recall}---detecting every planet with zero false negatives. The clustering approach proved effective for handling stellar diversity: by grouping windows with similar statistical properties (variability, noise characteristics), the model learned specialized representations for different stellar types. Hyperparameter optimization via Optuna identified optimal settings of 4 BiLSTM layers, 192 hidden units, 0.334 dropout, and batch size 128. The 50$\times$ dataset expansion from 655 to 33,051 windows enabled training deeper networks without overfitting, resulting in a \textbf{33.3\% improvement} over the baseline model. However, cross-mission generalization remains a significant limitation: fine-tuning on mixed TESS+Kepler data actually degraded Kepler detection from 3.7\% to 0\%, demonstrating that the model learns mission-specific features rather than universal transit physics.

\textbf{CNN Findings.} [TO BE COMPLETED BY TRISTAN]

\textbf{Transformer Findings.} [TO BE COMPLETED BY BREE]

\textbf{Comparative Analysis.} [TO BE COMPLETED AFTER ALL RESULTS AVAILABLE - will compare AUC, precision, recall, and computational cost across all three architectures]

\textbf{Contributions.} This work makes three contributions: (1) a systematic comparison of CNN, RNN, and Transformer architectures on identical TESS data; (2) demonstration that unsupervised clustering of statistical features can improve transit detection by enabling model specialization; and (3) open-source implementations of all three approaches for reproducibility.

\section*{AI Use Disclosure}

Portions of this work were completed with assistance from \textbf{Claude Code}, using models \texttt{claude-sonnet-4-5-20250929} (Sonnet 4.5) and \texttt{claude-opus-4-5-20251101} (Opus 4.5).

AI assistance for the RNN component included:
\begin{itemize}
\item Deeper understanding RNN concepts from presentation feedback and research papers
\item Assisted in implementing BiLSTM architecture to work with PyTorch
\item Debugging data pipeline and training code (handling Windows multiprocessing issues, GPU memory management)
\item Writing comprehensive documentation and comments in code
\item Experiment design and hyperparameter tuning guidance
\item LaTeX formatting and table generation
\end{itemize}

AI assistance for the CNN component included:
\begin{itemize}
\item Deeper understanding of CNN components from papers
\item Assisted in giving different approaches to are developed CNN architecture
\item Debugging data pipeline and preprocessing
\end{itemize}

All code was reviewed, understood, and tested by the students before submission. All experimental results were obtained by running code on real data. All design decisions (architecture choices, hyperparameters, evaluation metrics) were made by the students. The AI did not complete the project autonomously -- student understanding, testing, debugging, and decision-making were central to all work.


\bibliography{resourceFile}

\end{document}
